\documentclass[11pt]{article}
\usepackage[margin=1in]{geometry}
\usepackage{amsmath,amssymb,amsthm,mathtools,microtype,graphicx}
\usepackage[hidelinks]{hyperref}
\newtheorem{theorem}{Theorem}
\newtheorem{lemma}{Lemma}
\newcommand{\R}{\mathbb R}
\newcommand{\conv}{\operatorname{conv}}
\newcommand{\Span}{\operatorname{span}}
\setlength{\parindent}{0pt}
\setlength{\parskip}{0.62em}

\title{A Unit-Norm Full-Spark Frame That Is Not Splittable\\
\large Counterexample to the full-spark splittability conjecture}
\author{Run \texttt{fa\_banach\_001}, lane 18}
\date{11 August 2026}

\begin{document}
\maketitle

\begin{center}
\fbox{\parbox{0.91\textwidth}{\textbf{Status.}
Candidate full counterexample, likely valid, for human review. The proof is
exact and finite-dimensional; novelty remains subject to broader literature
review.}}
\end{center}

\section{Source conjecture}

For a finite frame $F=(f_1,\ldots,f_n)$ spanning $\R^d$, define
\[
 \|z\|_F=\min\{\|q\|_1:Fq=z\},
 \qquad
 \Sigma_{F,s}=\bigcup_{|S|\leq s}\Span\{f_j:j\in S\}.
\]
Choose a best $s$-term approximant
\[
 z_s\in\mathop{\rm argmin}_{v\in\Sigma_{F,s}}\|z-v\|_F.
\]
Chen, Kummerle, and Wang call $F$ \emph{$s$-splittable} with constant
$\beta>0$ if, for every $x,y\in\R^d$,
\begin{equation}\label{eq:split}
 \|x+y\|_F\geq \|x_s\|_F-\|y_s\|_F
 +\beta\bigl(\|y-y_s\|_F-\|x-x_s\|_F\bigr).
\end{equation}
Immediately after the definition they conjecture that every full-spark frame
is splittable for a sufficiently small positive $\beta$ \cite{CKW}. The
construction below disproves this already for $d=2$ and $s=1$.

\begin{center}
 \includegraphics[width=\textwidth]{figures/open_problem_crop.png}

 \small Definition 3.6 and the full-spark splittability conjecture,
 source PDF page 7.
\end{center}

\section{Idea of the proof}

With three atoms in the plane, the atomic unit ball is a centrally symmetric
hexagon, so both its norm and the distance to each atom line have exact
piecewise-linear formulas. Every pair $(x,y)$ imposes a linear constraint
$L\geq\beta D$ on the purported splitting constant. We choose rational unit
atoms and two rational witness pairs: the first has $L,D<0$ and forces a
strict lower bound on $\beta$, while the second has $L,D>0$ and forces a
smaller upper bound. Strict distance comparisons and exposed facets make all
four best one-term approximants unique.

\section{Counterexample and exact atomic norm}

Set
\begin{equation}\label{eq:frame}
 a=\left(-\frac7{25},-\frac{24}{25}\right),\qquad
 b=(1,0),\qquad
 c=\left(\frac{20}{29},-\frac{21}{29}\right),
 \qquad F=(a,b,c).
\end{equation}

\begin{theorem}\label{thm:main}
The frame in \eqref{eq:frame} is a unit-norm full-spark frame for $\R^2$, but
it is not $1$-splittable for any positive constant $\beta$. In fact, no real
$\beta$ satisfies \eqref{eq:split} for all $x,y$.
\end{theorem}

Each column has Euclidean norm one, and
\[
 \det(a,b)=\frac{24}{25},\qquad
 \det(a,c)=\frac{627}{725},\qquad
 \det(b,c)=-\frac{21}{29}.
\]
Thus every pair is linearly independent, so $F$ is full spark.

Write $N(z)=\|z\|_F$. Its unit ball is the centrally symmetric hexagon
$\conv\{\pm a,\pm b,\pm c\}$. The three pairs of opposite supporting lines
give the exact formula
\begin{equation}\label{eq:norm}
 N(z_1,z_2)=\max\left\{
 \left|z_1+\frac34z_2\right|,
 \left|-\frac3{11}z_1+\frac{37}{33}z_2\right|,
 \left|-z_1+\frac37z_2\right|
 \right\}.
\end{equation}
Indeed, the first functional takes the values $-1,1,17/116$ on $a,b,c$;
the second takes $-1,-3/11,-1$; and the third takes $-23/175,-1,-1$.
Hence each has absolute value at most one on the six vertices and exposes the
claimed facet. These are all three pairs of facets of the hexagon, proving
\eqref{eq:norm}.

\section{Best one-term approximation}

For $v\in\{a,b,c\}$, let
\[
 d_v(z)=\min_{t\in\R}N(z-tv).
\]

\begin{lemma}\label{lem:distances}
The distances to the three atom lines are
\begin{align}
 d_a(z)&=\frac{|24z_1-7z_2|}{24},\label{eq:da}\\
 d_b(z)&=\frac{25|z_2|}{24},\label{eq:db}\\
 d_c(z)&=\frac{25|21z_1+20z_2|}{627}.\label{eq:dc}
\end{align}
Whenever one of these values is strictly smaller than the other two, the best
one-term approximant is unique.
\end{lemma}

\begin{proof}
The annihilator $(24,-7)$ of $a$ has dual norm
\[
 \max\{|(24,-7)a|,|(24,-7)b|,|(24,-7)c|\}=24,
\]
with the maximum attained only at $\pm b$. The quotient norm modulo
$\Span(a)$ is therefore \eqref{eq:da}. Equality can occur only when the
residual is a multiple of $b$, which determines the minimizing point on
$\Span(a)$ uniquely. Similarly, $(0,1)$ annihilates $b$, has dual norm
$24/25$ attained only at $\pm a$, and gives \eqref{eq:db}. Finally,
$(21,20)$ annihilates $c$, has dual norm $627/25$ attained only at $\pm a$,
and gives \eqref{eq:dc}. Thus the minimizing point on each selected line is
unique. Strict inequality between the three distances also makes the selected
atom line unique.
\end{proof}

\section{Two incompatible constraints on the splitting constant}

Take first
\[
 x=\left(-\frac45,\frac35\right),\qquad
 y=\left(\frac7{15},-\frac8{15}\right).
\]
Lemma~\ref{lem:distances} gives
\[
 (d_a(x),d_b(x),d_c(x))
 =\left(\frac{39}{40},\frac58,\frac{40}{209}\right),
\]
and
\[
 (d_a(y),d_b(y),d_c(y))
 =\left(\frac{28}{45},\frac59,\frac{65}{1881}\right).
\]
Thus both best one-term approximants are unique and lie on $\Span(c)$. They
and their residuals are
\begin{align*}
 x_1&=-\frac{1131}{1045}c
 =\left(-\frac{156}{209},\frac{819}{1045}\right),
 &x-x_1&=\frac{40}{209}a,\\
 y_1&=\frac{6496}{9405}c
 =\left(\frac{896}{1881},-\frac{1568}{3135}\right),
 &y-y_1&=\frac{65}{1881}a.
\end{align*}
Formula \eqref{eq:norm} now yields
\[
 N(x_1)=\frac{1131}{1045},\quad
 N(y_1)=\frac{6496}{9405},\quad
 N(x+y)=N\left(-\frac13,\frac1{15}\right)=\frac{38}{105}.
\]
After moving the non-$\beta$ terms in \eqref{eq:split} to the left, this
witness gives
\begin{equation}\label{eq:lower}
 \underbrace{N(x+y)-N(x_1)+N(y_1)}_{-391/13167}
 \geq
 \beta\underbrace{\bigl(N(y-y_1)-N(x-x_1)\bigr)}_{-295/1881}.
\end{equation}
Both quantities are negative, so \eqref{eq:lower} forces
\begin{equation}\label{eq:lowerbound}
 \beta\geq\frac{391}{2065}.
\end{equation}

\clearpage
For the second witness, take
\[
 x'=b=(1,0),\qquad y'=\left(-\frac3{10},-\frac15\right).
\]
The best approximant to $x'$ is uniquely $x'_1=b$. For $y'$ the three
distances are
\[
 (d_a(y'),d_b(y'),d_c(y'))
 =\left(\frac{29}{120},\frac5{24},\frac{515}{1254}\right),
\]
so its unique best approximant lies on $\Span(b)$, with
\[
 y'_1=\left(-\frac{29}{120},0\right),\qquad
 y'-y'_1=\frac5{24}a.
\]
Moreover
\[
 N(x'+y')=N\left(\frac7{10},-\frac15\right)=\frac{11}{14}.
\]
Consequently \eqref{eq:split} gives
\begin{equation}\label{eq:upper}
 \underbrace{\frac{11}{14}-1+\frac{29}{120}}_{23/840}
 \geq \beta\underbrace{\frac5{24}}_{N(y'-y'_1)-N(x'-x'_1)},
\end{equation}
which forces
\begin{equation}\label{eq:upperbound}
 \beta\leq\frac{23}{175}.
\end{equation}
But
\[
 \frac{391}{2065}-\frac{23}{175}=\frac{598}{10325}>0.
\]
The necessary bounds \eqref{eq:lowerbound} and \eqref{eq:upperbound} are
incompatible. This proves Theorem~\ref{thm:main}.

\section{Verification, scope, and novelty status}

All displayed values are rational. The accompanying standard-library script
\texttt{verify\_exact.py} checks the minors, facet inequalities,
distance comparisons, decompositions, norms, and final gap using exact
rational arithmetic. It is an audit aid; the proof above is self-contained.

The general definition can involve nonunique best sparse approximants
\cite{Chen}. Every best approximant used here is unique, so the counterexample
does not depend on how an $\mathop{\rm argmin}$ is selected. Unit-norm columns
also show that the obstruction is not caused by rescaling a frame vector.

A bounded search of the run indexes and exact web queries found the source
conjecture and the earlier atomic-norm article, but no later proof or
counterexample. Priority is not asserted. The result should receive expert
proof review and a broader citation search before dissemination.

\begin{thebibliography}{9}
\bibitem{CKW}
X.~Chen, C.~Kummerle, and R.~Wang,
\emph{Sparse Recovery for Overcomplete Frames: Sensing Matrices and Recovery
Guarantees}, arXiv:2408.16166 (2024).

\bibitem{Chen}
X.~Chen,
\emph{A Unified Recovery of Structured Signals Using Atomic Norm},
Information and Inference 13 (2024), iaad050; arXiv:2207.06484.
\end{thebibliography}

\end{document}
